
\setcounter{footnote}{10}

\SGARefTarget{sga3:viii:section:2}\section*{2. Scheme-theoretic Properties of Diagonalizable Groups}
\addcontentsline{toc}{section}{2. Scheme-theoretic Properties of Diagonalizable Groups}

These properties are summarized in \SGARefLink{sga3:viii:proposition:2:1}{the following proposition}.

\medskip
\noindent\SGARefTarget{sga3:viii:proposition:2:1}\textbf{Proposition 2.1.}
\textit{Let \(S\) be a nonempty prescheme, let \(M\) be an ordinary
commutative group, and let \(G=D(M_S)\) be the diagonalizable \(S\)-group
defined by \(M\). Then:}
\begin{enumerate}
  \item[(a)]\SGARefTarget{sga3:viii:proposition:2:1:item:a:01}
  \textit{\(G\) is faithfully flat over \(S\) and affine over \(S\)
  (a fortiori, quasi-compact over \(S\));}

  \item[(b)]\SGARefTarget{sga3:viii:proposition:2:1:item:b:01}
  \textit{\(M\) is of finite type if and only if \(G\) is of finite type
  over \(S\), if and only if \(G\) is of finite presentation over \(S\);}

  \item[(c)]\SGARefTarget{sga3:viii:proposition:2:1:item:c:01}
  \textit{\(M\) is finite if and only if \(G\) is finite over \(S\), if
  and only if \(G\) is of finite type over \(S\) and is annihilated by an
  integer \(n>0\). In that case,}
  \[
    \deg(G/S)=\operatorname{Card}(M);
  \]

  \item[\((c')\)]\SGARefTarget{sga3:viii:proposition:2:1:item:cprime:01}
  \textit{\(M\) is a torsion group if and only if \(G\) is integral over
  \(S\);}

  \item[(d)]\SGARefTarget{sga3:viii:proposition:2:1:item:d:01}
  \textit{\(M=0\) if and only if \(G\) is the unit \(S\)-group;}

  \item[(e)]\SGARefTarget{sga3:viii:proposition:2:1:item:e:01}
  \textit{\(M\) is of finite type and the order of its torsion subgroup is
  prime to the residue characteristics of \(S\) if and only if \(G\) is
  smooth over \(S\).}
\end{enumerate}

The verification of \SGARefLink{sga3:viii:proposition:2:1:item:a:01}{(a)}
through \SGARefLink{sga3:viii:proposition:2:1:item:d:01}{(d)} is immediate
and is left to the reader. Let us prove
\SGARefLink{sga3:viii:proposition:2:1:item:e:01}{(e)}. If \(G\) is smooth
over \(S\), it is locally of finite presentation over \(S\), hence of
finite presentation over \(S\), since it is affine over \(S\); thus \(M\)
is of finite type. We may therefore assume from the outset that \(M\) is
of finite type, and hence that \(G\) is of finite presentation over
\(S\). Then\footnote{Editor's note: this uses the fact that \(G\) is flat
over \(S\), by \SGARefLink{sga3:viii:proposition:2:1:item:a:01}{(a)}.}
\(G\) is smooth over \(S\) if and only if its geometric fibers are
smooth, reducing us to the case in which \(S\) is the spectrum of an
algebraically closed field \(k\).

Write \(M=T\times L\), where \(T\) is the torsion subgroup and \(L\) is
free, with \(L\simeq\mathbb Z^r\). Then
\[
  D(M)=D(T)\times D(L),
  \qquad
  D(L)\simeq\mathbb G_m^r,
\]
and \(D(L)\) is smooth over \(k\). Thus \(G=D(M)\) is smooth over \(k\)
if and only if \(D(T)\) is smooth. Since \(D(T)\) is finite over \(k\) of
degree equal to the order \(n\) of \(T\), this means that \(D(T)(k)\) has
\(n\) elements. Now \(T\) is isomorphic to a sum of groups
\(\mathbb Z/n_i\mathbb Z\), where \(n=\prod_i n_i\); consequently
\[
  D(T)=\prod_iD(\mathbb Z/n_i\mathbb Z)
      =\prod_i\mu_{n_i},
\]
where \(\mu_{n_i}\) is the group scheme of \(n_i\)-th roots of unity.
Hence
\[
  \operatorname{card}\bigl(D(T)(k)\bigr)
  =\prod_i\operatorname{card}\bigl(\mu_{n_i}(k)\bigr).
\]
Here \(\operatorname{card}(\mu_{n_i}(k))\), the number of \(n_i\)-th
roots of unity in \(k\), is at most \(n_i\), and equality holds if and
only if \(n_i\) is prime to the characteristic \(p\) of \(k\).
\emph{[The source prints \(\mu_i(k)\) in both occurrences.]}
It follows
that
\[
  \operatorname{card}\bigl(D(T)(k)\bigr)=n,
  \qquad n=\prod_i n_i,
\]
if and only if every \(n_i\) is prime to \(p\), that is, if and only if
\(n\) is prime to \(p\).
\hfill\(\square\)

\SGARefTarget{sga3:viii:section:3}\section*{3. Exactness Properties of the Functor
\texorpdfstring{\(D_S\)}{D\_S}}
\addcontentsline{toc}{section}{3. Exactness Properties of the Functor
\texorpdfstring{\(D_S\)}{D\_S}}

\medskip
\noindent\SGARefTarget{sga3:viii:theorem:3:1}\textbf{Theorem 3.1.}
\textit{Let \(S\) be a prescheme, and let}
\SGARefTarget{sga3:viii:theorem:3:1:diagram:01}
\[
  0\longrightarrow M'
   \xrightarrow{u}M
   \xrightarrow{v}M''
   \longrightarrow0
\]
\textit{be an exact sequence of ordinary commutative groups. Consider the
sequence of transposed homomorphisms}
\SGARefTarget{sga3:viii:theorem:3:1:diagram:02}
\[
  0\longrightarrow D_S(M'')
   \xrightarrow{v^t}D_S(M)
   \xrightarrow{u^t}D_S(M')
   \longrightarrow0.
\]
\begin{enumerate}
  \item[(i)]\SGARefTarget{sga3:viii:theorem:3:1:item:i:01}
  \textit{The homomorphism \(v^t\) induces an isomorphism from
  \(D_S(M'')\) onto the kernel of \(u^t\), and \(u^t\) is faithfully flat
  and quasi-compact.}

  \item[(ii)]\SGARefTarget{sga3:viii:theorem:3:1:item:ii:01}%
  \footnote{Editor's note: what follows has been added, and the proof has
  been expanded accordingly.}
  \textit{\(D_S(M')\) represents the quotient sheaf
  \(D_S(M)/D_S(M'')\) for the \((\mathrm{fpqc})\) topology.}
\end{enumerate}

Let \(\mathcal M\) denote the family of faithfully flat quasi-compact
morphisms. First,
\SGARefLink{sga3:viii:theorem:3:1:item:ii:01}{(ii)} follows from
\SGARefLink{sga3:viii:theorem:3:1:item:i:01}{(i)} (cf. IV,
\SGARefLink{sga3:iv:corollary:4:6:5:1}{4.6.5.1}). Indeed, the equivalence
relation on \(D_S(M)\) defined by \(u^t\) is the same as that defined by
the subgroup
\[
  \operatorname{Ker}(u^t)=D_S(M'').
\]
Since \(u^t\in\mathcal M\), this equivalence relation is
\(\mathcal M\)-effective (cf. IV,
\SGARefLink{sga3:iv:scholium:3:3:2:1}{3.3.2.1}); therefore \(D_S(M')\)
represents the quotient sheaf for the \((\mathrm{fpqc})\) topology (cf.
IV, \SGARefLink{sga3:iv:lemma:4:6:5}{4.6.5}).

The first assertion of
\SGARefLink{sga3:viii:theorem:3:1:item:i:01}{(i)} is an immediate
consequence of the definition of the functors \(D_S(-)\). More generally,
every exact sequence
\SGARefTarget{sga3:viii:theorem:3:1:diagram:03}
\[
  M'\longrightarrow M\longrightarrow M''\longrightarrow0
\]
with no zero on the left gives rise to an exact transposed sequence
\SGARefTarget{sga3:viii:theorem:3:1:diagram:04}
\[
  0\longrightarrow D_S(M'')
   \longrightarrow D_S(M)
   \longrightarrow D_S(M').
\]
(This remains valid in the more general setting considered at the
beginning of \(\SGARefLink{sga3:viii:section:1}{\S1}\).) On the other hand, because \(D_S(M)\) and
\(D_S(M')\) are affine over \(S\), the morphism \(u^t\) is necessarily
affine, and a fortiori quasi-compact, whatever the homomorphism
\(u:M'\to M\). The second assertion of
\SGARefLink{sga3:viii:theorem:3:1:item:i:01}{(i)} will therefore follow
from \SGARefLink{sga3:viii:corollary:3:2:item:a:01}{part (a)} of \SGARefLink{sga3:viii:corollary:3:2}{the
following corollary}.

\medskip
\noindent\SGARefTarget{sga3:viii:corollary:3:2}\textbf{Corollary 3.2.}
\textit{Let \(S\) be a nonempty prescheme, let \(u:M'\to M\) be a
homomorphism of ordinary commutative groups, and let
\(u^t:G\to G'\) be the transposed homomorphism. Then:}
\begin{enumerate}
  \item[(a)]\SGARefTarget{sga3:viii:corollary:3:2:item:a:01}
  \textit{\(u\) is a monomorphism if and only if \(u^t\) is faithfully
  flat;}

  \item[(b)]\SGARefTarget{sga3:viii:corollary:3:2:item:b:01}
  \textit{\(u\) is an epimorphism if and only if \(u^t\) is a
  monomorphism; in that case \(u^t\) is in fact a closed immersion.}
\end{enumerate}

To prove \SGARefLink{sga3:viii:corollary:3:2:item:a:01}{(a)}, observe
that if \(u\) is a monomorphism, then \(\mathcal O_S(M)\), as a module
over \(\mathcal O_S(M')\), has a nonempty basis: namely, the system of
sections defined by any system of representatives of \(M\) modulo
\(M'\). It is therefore faithfully flat. Conversely, if
\(\mathcal O_S(M)\) is faithfully flat over \(\mathcal O_S(M')\), then
\[
  u^t:\mathcal O_S(M')\longrightarrow\mathcal O_S(M)
\]
is injective, which, when \(S\ne\varnothing\), implies that
\(u:M'\to M\) is injective.

To prove \SGARefLink{sga3:viii:corollary:3:2:item:b:01}{(b)}, observe
that if \(u\) is an epimorphism, then
\(\mathcal O_S(M')\to\mathcal O_S(M)\) is surjective, so \(u^t\) is a
closed immersion and, a fortiori, a monomorphism. Conversely, if \(u^t\)
is a monomorphism, then \(\operatorname{Ker}u^t\) is the unit group.
Putting \(M''=\operatorname{Coker}u\), we have seen that
\[
  \operatorname{Ker}u^t\simeq D_S(M'').
\]
Thus \SGARefLink{sga3:viii:proposition:2:1:item:d:01}{2.1(d)} gives
\(M''=0\), and \(u\) is an epimorphism.

The usual argument now gives \SGARefLink{sga3:viii:corollary:3:3}{the following consequence} of
\SGARefLink{sga3:viii:theorem:3:1}{3.1}.

\medskip
\noindent\SGARefTarget{sga3:viii:corollary:3:3}\textbf{Corollary 3.3.}
\textit{Let}
\SGARefTarget{sga3:viii:corollary:3:3:diagram:00}
\[
  M'\xrightarrow{u}M\xrightarrow{v}M''
\]
\textit{be an exact sequence of ordinary commutative groups, and consider
the transposed sequence}
\SGARefTarget{sga3:viii:corollary:3:3:diagram:01}
\[
  G''\xrightarrow{v^t}G\xrightarrow{u^t}G'.
\]
\textit{Then \(v^t\) induces a faithfully flat quasi-compact morphism
from \(G''\) to \(\operatorname{Ker}u^t\), and this kernel is a
diagonalizable group isomorphic to}
\[
  D_S(v(M))=D_S(\operatorname{Coker}u).
\]

\medskip
\noindent\SGARefTarget{sga3:viii:corollary:3:4}\textbf{Corollary 3.4.}
\textit{Let \(S\) be a prescheme, let \(u:G\to H\) be a homomorphism of
\(S\)-group preschemes that are locally diagonalizable, with \(H\) of
finite type over \(S\), and put \(G'=\operatorname{Ker}u\). Then:}
\begin{enumerate}
  \item[(a)]\SGARefTarget{sga3:viii:corollary:3:4:item:a:01}
  \textit{\(G'\) is locally diagonalizable, and is of finite type over
  \(S\) if \(G\) is;}

  \item[(b)]\SGARefTarget{sga3:viii:corollary:3:4:item:b:01}
  \textit{the quotient \(G/G'\) ``exists'': more precisely, the
  equivalence relation on \(G\) defined by \(G'\) is
  \(\mathcal M\)-effective, where \(\mathcal M\) is the family of
  faithfully flat quasi-compact morphisms (cf. IV,
  \SGARefLink{sga3:iv:subsection:3:4}{3.4}). Moreover, \(G/G'\) is
  locally diagonalizable and of finite type over \(S\);}

  \item[(c)]\SGARefTarget{sga3:viii:corollary:3:4:item:c:01}
  \textit{the homomorphism \(u:G\to H\) factors uniquely as}
  \SGARefTarget{sga3:viii:corollary:3:4:diagram:01}
  \[
    G\xrightarrow{v}G/G'\xrightarrow{w}H,
  \]
  \textit{where \(v\) is the canonical homomorphism, and is therefore
  faithfully flat and quasi-compact. Moreover, \(w\) is a closed
  immersion, and a fortiori a monomorphism.}
\end{enumerate}
\textit{Finally, the quotient}
\[
  H'=H/\operatorname{Im}w
    =\operatorname{Coker}w
    =\operatorname{Coker}u
\]
\textit{exists: more precisely, the equivalence relation on \(H\)
defined by \(G/G'\) is \(\mathcal M\)-effective, and \(H'\) is of finite
type over \(S\).}

The first assertion of
\SGARefLink{sga3:viii:corollary:3:4:item:c:01}{(c)} follows from
\SGARefLink{sga3:viii:corollary:3:4:item:b:01}{(b)}, by the definition
of the quotient \(G/G'\) (cf. IV,
\SGARefLink{sga3:iv:lemma:3:2:3}{3.2.3}).
\footnote{Editor's note: the original has been expanded in what follows.}
Let us show that the \((\mathrm{fpqc})\) quotient sheaf
\(\widetilde G/\widetilde{G'}\) is representable. This can be checked
locally on \(S\), as can all the other assertions; we may therefore
suppose that \(G\) and \(H\) are diagonalizable, of the form \(D_S(M)\)
and \(D_S(N)\).

Since \(H\) is of finite type over \(S\), the group \(N\) is of finite
type by \SGARefLink{sga3:viii:proposition:2:1:item:b:01}{2.1(b)}.
Therefore, by \SGARefLink{sga3:viii:corollary:1:5}{1.5}, \(u\) is defined by a homomorphism \(u':N\to M\).
It follows from \SGARefLink{sga3:viii:theorem:3:1}{3.1} and
\SGARefLink{sga3:viii:corollary:3:2}{3.2} that
\[
  G'\simeq D_S(\operatorname{Coker}u')
\]
and that \(\widetilde G/\widetilde{G'}\) is represented by
\(D_S(\operatorname{Im}u')\). Moreover, on considering the exact
sequence
\SGARefTarget{sga3:viii:corollary:3:4:diagram:02}
\[
  0\longrightarrow\operatorname{Ker}u'
   \longrightarrow N
   \xrightarrow{w'}\operatorname{Im}u'
   \longrightarrow0,
\]
we find that \(w\) is a closed immersion and that
\[
  H'=H/\operatorname{Im}w=D_S(\operatorname{Ker}u').
\]
The latter is of finite type over \(S\), since \(N\), and hence
\(\operatorname{Ker}u'\), is of finite type.

\medskip
\noindent\SGARefTarget{sga3:viii:remarks:3:4}\textbf{Remarks.}
The existence result for quotients in
\SGARefLink{sga3:viii:corollary:3:4}{3.4} will be substantially
generalized in \(\SGARefLink{sga3:viii:section:5}{\S5}\).

On the other hand, note that in \SGARefLink{sga3:viii:section:3}{the present section} and \SGARefLink{sga3:viii:section:2}{the preceding
one}, the hypothesis \(S\ne\varnothing\) was used only to ensure the
validity of certain converses, allowing properties of the corresponding
ordinary groups to be deduced from certain hypotheses on diagonalizable
\(S\)-groups. The results in the ``forward'' direction hold without any
restriction on \(S\), and the proofs given here apply in general.

\medskip
\noindent\SGARefTarget{sga3:viii:corollary:3:5}\textbf{Corollary 3.5.}
\textit{Let \(G\) be a diagonalizable group prescheme over \(S\), and let
\(n\ne0\) be an integer. Then the subgroup \({}_nG\) of \(G\), the
kernel of the homomorphism}
\[
  n\mathbin{\cdot}\operatorname{id}_G:G\longrightarrow G,
\]
\textit{is integral over \(S\), and is finite over \(S\) if \(G\) is of
finite type over \(S\).}

Indeed, if \(G=D_S(M)\), then
\[
  {}_nG=D_S(M/nM)
\]
by \SGARefLink{sga3:viii:theorem:3:1}{3.1}, and the result follows from
\SGARefLink{sga3:viii:proposition:2:1:item:b:01}{2.1(b)},
\SGARefLink{sga3:viii:proposition:2:1:item:c:01}{(c)}, and
\SGARefLink{sga3:viii:proposition:2:1:item:cprime:01}{\((c')\)}.

\SGARefTarget{sga3:viii:section:4}\section*{4. Torsors under a Diagonalizable Group}
\addcontentsline{toc}{section}{4. Torsors under a Diagonalizable Group}

\SGARefTarget{sga3:viii:section:4:problem:01}Let \(S\) be a prescheme and let \(G=D_S(M)\) be a diagonalizable group
over \(S\). We shall determine the \(G\)-torsors (or principal
homogeneous \(G\)-bundles) over \(S\) for the faithfully flat
quasi-compact topology; cf. Expos\'e IV,
\SGARefLink{sga3:iv:subsection:5:1}{5.1}. Recall that a prescheme \(P\)
over \(S\), with \(S\)-operator group \(G\), is called a \emph{torsor} or
a \emph{principal homogeneous bundle} if every point of \(S\) has an
open neighborhood \(U\) and a faithfully flat quasi-compact morphism
\(S'\to U\) such that
\[
  P'=P\times_SS'
\]
is isomorphic, as an object with operators, to
\[
  G'=G\times_SS',
\]
where \(G'\) acts on itself by right translations. Since \(G\) is affine
over \(S\), SGA 1, VIII 5.6 implies that \(P\) is necessarily affine over
\(S\). Also, because \(G\) itself is faithfully flat and quasi-compact
over \(S\), \(P\) is principal homogeneous under \(G\) if and only if it
is formally principal homogeneous and is, in addition, faithfully flat
and quasi-compact over \(S\) (cf. IV,
\SGARefLink{sga3:iv:proposition:5:1:6}{5.1.6}).

Recall also (\SGARefLink{sga3:i:proposition:4-7-3}{Expos\'e I, 4.7.3}) that giving an \(S\)-prescheme \(P\), affine
over \(S\), with \(S\)-operator group \(G=D_S(M)\), amounts to giving a
commutative quasi-coherent algebra graded by \(M\) over \(S\). In other
words, it amounts to giving a quasi-coherent algebra \(\mathscr A\) over
\(S\), together with a direct-sum decomposition as a module
\[
  \mathscr A=\coprod_{m\in M}\mathscr A_m,
\]
such that
\[
  \mathscr A_m\mathbin{\cdot}\mathscr A_{m'}
  \subseteq\mathscr A_{m+m'}
  \qquad\text{for }m,m'\in M.
\]
With this understood, the answer to \SGARefLink{sga3:viii:section:4:problem:01}{the preceding problem} is as follows.

\medskip
\noindent\SGARefTarget{sga3:viii:proposition:4:1}\textbf{Proposition 4.1.}
\textit{Let \(P\) be the prescheme with \(S\)-operator group
\(G=D_S(M)\) defined by the \(M\)-graded algebra \(\mathscr A\). Then
\(P\) is a principal homogeneous bundle under \(G\) if and only if
\(\mathscr A\) satisfies the following conditions:}
\begin{enumerate}
  \item[(a)]\SGARefTarget{sga3:viii:proposition:4:1:item:a:01}
  \textit{for every \(m\in M\), \(\mathscr A_m\) is an invertible module
  over \(S\);}

  \item[(b)]\SGARefTarget{sga3:viii:proposition:4:1:item:b:01}
  \textit{for \(m,m'\in M\), the homomorphism}
  \[
    \mathscr A_m\otimes_{\mathcal O_S}\mathscr A_{m'}
    \longrightarrow\mathscr A_{m+m'}
  \]
  \textit{induced by multiplication in \(\mathscr A\) is an
  isomorphism.}
\end{enumerate}

The necessity of \SGARefLink{sga3:viii:proposition:4:1}{these conditions} follows immediately by descent, since
they hold when \(P\) is the trivial principal homogeneous bundle, that
is, when \(\mathscr A=\mathcal O_S(M)\). For sufficiency, note that
\SGARefLink{sga3:viii:proposition:4:1:item:a:01}{(a)} already implies
that \(P\) is faithfully flat over \(S\); it is in any event
quasi-compact over \(S\), since it is affine over \(S\). It therefore
remains to check that \(P\) is formally principal homogeneous under
\(G\), that is, that the familiar morphism
\SGARefTarget{sga3:viii:proposition:4:1:diagram:01}
\[
  P\times_SG\longrightarrow P\times_SP
\]
is an isomorphism. On affine algebras \SGARefLink{sga3:viii:proposition:4:1:diagram:01}{this morphism} is given explicitly
by the homomorphism
\SGARefTarget{sga3:viii:proposition:4:1:diagram:02}
\[
  \mathscr A\otimes\mathscr A
  \longrightarrow
  \mathscr A(M)=\mathscr A\otimes\mathcal O_S(M).
\]

In bidegree \((m,n)\), where \(m,n\in M\), it is given by
\[
  x_m\otimes y_n\longmapsto x_my_n\otimes e_n.
\]
At the level of degrees, \SGARefLink{sga3:viii:proposition:4:1:diagram:02}{this homomorphism} is compatible with the
homomorphism \(M\times M\to M\times M\) given by
\[
  (m,n)\longmapsto(m+n,n),
\]
which is an isomorphism. This shows that
\SGARefLink{sga3:viii:proposition:4:1:item:b:01}{(b)} expresses precisely,
and independently of \SGARefLink{sga3:viii:proposition:4:1:item:a:01}{(a)},
that \(P\) is formally principal homogeneous, and proves
\SGARefLink{sga3:viii:proposition:4:1}{4.1}.

Faithfully flat descent also gives:

\medskip
\noindent\SGARefTarget{sga3:viii:corollary:4:2}\textbf{Corollary 4.2.}
\textit{The conditions of \SGARefLink{sga3:viii:proposition:4:1}{4.1}
imply that the homomorphism}
\[
  \mathcal O_S\longrightarrow\mathscr A_0
\]
\textit{is an isomorphism.}

For example, if \(M=\mathbb Z\), then, under the conditions of
\SGARefLink{sga3:viii:proposition:4:1}{4.1}, the algebra \(\mathscr A\)
is essentially determined by \(\mathscr A_1=\mathscr L\), namely
\[
  \mathscr A\simeq\coprod_{n\in\mathbb Z}\mathscr L^{\otimes n}
\]
as graded algebras. Thus we recover the familiar result:

\medskip
\noindent\SGARefTarget{sga3:viii:corollary:4:3}\textbf{Corollary 4.3.}
\textit{The category of principal homogeneous bundles \(P\) over \(S\)
under \(\mathbb G_{m,S}\) is equivalent to the category of invertible
modules \(\mathscr L\) over \(S\), where the morphisms in both categories
are the isomorphisms of the structures in question. Quasi-inverse
functors are obtained by assigning to \(P\) the degree-\(1\) component of
its \(\mathbb Z\)-graded affine algebra, and by assigning to
\(\mathscr L\) the spectrum of the \(\mathbb Z\)-graded algebra}
\[
  \coprod_{n\in\mathbb Z}\mathscr L^{\otimes n}.
\]

In particular:

\medskip
\noindent\SGARefTarget{sga3:viii:corollary:4:4}\textbf{Corollary 4.4.}
\textit{The group of isomorphism classes of principal homogeneous
bundles over \(S\) under \(\mathbb G_{m,S}\) is isomorphic to the group
\(\operatorname{Pic}(S)\) of isomorphism classes of invertible modules
over \(S\), that is, to}
\[
  H^1(S,\mathcal O_S^\times).
\]

Since \(\mathbb G_{m,S}\) is the scheme of automorphisms of the module
\(\mathcal O_S\), Corollary \SGARefLink{sga3:viii:corollary:4:4}{4.4} is
equivalent to \SGARefLink{sga3:viii:corollary:4:5}{the following statement}, one of the forms of Hilbert's
``\SGARefLink{sga3:viii:corollary:4:5}{Theorem 90}.''

\medskip
\noindent\SGARefTarget{sga3:viii:corollary:4:5}\textbf{Corollary 4.5.}
\textit{Every principal homogeneous bundle over \(S\) under
\(\mathbb G_m\) is locally trivial for the Zariski topology.}

\medskip
\noindent\SGARefTarget{sga3:viii:remark:4:5:1}\textbf{Remark 4.5.1.}
\SGARefLink{sga3:viii:corollary:4:5}{The preceding statement} is no longer true in general for a group such as
\(\mu_n\), or for a ``twisted form'' of \(\mathbb G_m\). For example,
the unique twisted form of \(\mathbb G_m\) over the field
\(\mathbb R\) of real numbers has a first cohomology group equal to
\(\mathbb Z/2\mathbb Z\).

\footnote{Editor's note: what follows has been added.}
Indeed, let \(\mathbb S^1\) be the kernel of the norm morphism
\[
  N:\prod_{\mathbb C/\mathbb R}\mathbb G_{m,\mathbb C}
  \longrightarrow\mathbb G_{m,\mathbb R};
\]
it is a \(\mathbb C/\mathbb R\)-twisted form of
\(\mathbb G_{m,\mathbb R}\). The equation \(N(z)=-1\) in
\(\prod_{\mathbb C/\mathbb R}(\mathbb G_{m,\mathbb C})\) defines an
\(\mathbb S^1\)-torsor \(X\) over \(\operatorname{Spec}(\mathbb R)\),
locally trivial for the \(\acute{e}\)tale topology, but not trivial, since
\(X(\mathbb R)=\varnothing\). Let us show that
\[
  H^1_{\mathrm{\acute et}}(\mathbb R,\mathbb S^1)
  \simeq\mathbb Z/2\mathbb Z.
\]
There is an exact sequence of smooth commutative algebraic
\(\mathbb R\)-groups
\SGARefTarget{sga3:viii:remark:4:5:1:diagram:01}
\[
  1\longrightarrow\mathbb S^1
  \longrightarrow\prod_{\mathbb C/\mathbb R}\mathbb G_{m,\mathbb C}
  \longrightarrow\mathbb G_{m,\mathbb R}
  \longrightarrow1,
\]
which gives rise to a long exact sequence in \(\acute{e}\)tale cohomology
(or in Galois cohomology):
\SGARefTarget{sga3:viii:remark:4:5:1:diagram:02}
\[
\begin{aligned}
  0\longrightarrow\mathbb S^1(\mathbb R)
  &\longrightarrow\mathbb C^\times
   \xrightarrow{N}\mathbb R^\times
   \longrightarrow H^1_{\mathrm{\acute et}}(\mathbb R,\mathbb S^1)\\
  &\longrightarrow
   H^1_{\mathrm{\acute et}}\!
   \left(\mathbb R,
     \prod_{\mathbb C/\mathbb R}\mathbb G_{m,\mathbb C}\right)
   \longrightarrow\cdots .
\end{aligned}
\]
Now (see, for example, XXIV, \SGARefLink{sga3:xxiv:proposition:8-4}{8.4}),
\[
  H^1_{\mathrm{\acute et}}\!
  \left(\mathbb R,
    \prod_{\mathbb C/\mathbb R}\mathbb G_{m,\mathbb C}\right)
  \simeq
  H^1_{\mathrm{\acute et}}(\mathbb C,\mathbb G_{m,\mathbb C}),
\]
and the latter group is zero by
\SGARefLink{sga3:viii:corollary:4:5}{4.5} (or, in this case, because
\(\mathbb C\) is algebraically closed). We therefore obtain an
isomorphism
\[
  H^1_{\mathrm{\acute et}}(\mathbb R,\mathbb S^1)
  \simeq\mathbb R^\times/N(\mathbb C^\times)
  \simeq\{\pm1\}.
\]

We shall need \SGARefLink{sga3:viii:proposition:4:6}{the following result} in \SGARefLink{sga3:viii:section:5}{the next section}.

\medskip
\noindent\SGARefTarget{sga3:viii:proposition:4:6}\textbf{Proposition 4.6.}
\textit{Under the conditions of
\SGARefLink{sga3:viii:proposition:4:1}{4.1}, conditions
\SGARefLink{sga3:viii:proposition:4:1:item:a:01}{(a)} and
\SGARefLink{sga3:viii:proposition:4:1:item:b:01}{(b)} are equivalent to
the following conditions:}
\begin{enumerate}
  \item[\((a')\)]\SGARefTarget{sga3:viii:proposition:4:6:item:aprime:01}
  \textit{the homomorphism
  \(\mathcal O_S\to\mathscr A_0\) is an isomorphism;}

  \item[\((b')\)]\SGARefTarget{sga3:viii:proposition:4:6:item:bprime:01}
  \textit{for every \(m\in M\) (it is enough to require this for the
  members of a system of generators of \(M\)), one has}
  \[
    \mathscr A_m\mathbin{\cdot}\mathscr A_{-m}=\mathscr A_0.
  \]
\end{enumerate}

The necessity is clear in view of
\SGARefLink{sga3:viii:corollary:4:2}{4.2}.\footnote{Editor's note:
``\SGARefLink{sga3:viii:proposition:4:1}{4.1}'' has been corrected to ``\SGARefLink{sga3:viii:corollary:4:2}{4.2}.''} We are thus reduced to proving
\SGARefLink{sga3:viii:corollary:4:7}{the following corollary}.

\medskip
\noindent\SGARefTarget{sga3:viii:corollary:4:7}\textbf{Corollary 4.7.}
\textit{Let}
\[
  A=\coprod_{n\in\mathbb Z}A_n
\]
\textit{be a \(\mathbb Z\)-graded ring such that}
\[
  A_1\mathbin{\cdot}A_{-1}=A_0.
\]
\textit{Then the \(A_n\) are invertible \(A_0\)-modules, and, for
\(n,n'\in\mathbb Z\), the homomorphism}
\[
  A_n\otimes_{A_0}A_{n'}\longrightarrow A_{n+n'}
\]
\textit{induced by multiplication in \(A\) is an isomorphism.}

By hypothesis there exist \(f_i\in A_1\) and \(g_i\in A_{-1}\) such that
\SGARefTarget{sga3:viii:corollary:4:7:eq:star:01}
\[
  \makebox[\displaywidth][s]{%
    \(\text{(*)}\)\hfill
    \(\displaystyle\sum_i f_ig_i=1.\)\hfill
    \(\phantom{\text{(*)}}\)%
  }
\]
Since the assertion to be proved is local on
\(\operatorname{Spec}(A_0)\),\footnote{Editor's note: the following
passage has been expanded.} and since
\(\SGARefLink{sga3:viii:corollary:4:7:eq:star:01}{(*)}\) shows that
\(\operatorname{Spec}(A_0)\) is covered by the affine open subsets
\(D(f_ig_i)\), we are reduced to the case in which there is an element
\(f\in A_1\) that is invertible in \(A\). Then, for every
\(n\in\mathbb Z\), \(f^n\) is an invertible element of \(A_n\), and
therefore defines an isomorphism
\[
  A_0\xrightarrow{\ \sim\ }A_n,
  \qquad h\longmapsto f^nh.
\]
Moreover,
